\begin{figure}[p]
\begin{subfigure}{\linewidth}
\centering
\includegraphics[trim={0 65mm 0 21.75mm}, clip, width=\linewidth, page=406]{images/screen2-texts}
\par\vspace{3.5pt}
\includegraphics[trim={0 7mm 0 75mm}, clip, width=\linewidth, page=406]{images/screen2-texts}
\par\vspace{-1.5pt}
\caption{Pro vrchní část panelu jsem použil řecký a anglický text z kanadské premiéry.~\cite{morland2018fantasia}}
\label{fig:screen2-texts-greek-and-english-original}
\end{subfigure}
\par\vspace{10pt}
\begin{subfigure}{\linewidth}
\begingroup
\footnotesize
\begin{verbatim}
$ pdftk dokument.pdf output dekomprimovaný-dokument.pdf uncompress
$ sed -i dekomprimovaný-dokument.pdf \
> -e's/\.686 \.933 \.933 RG \.686 \.933 \.933 rg/1 .890 .902 RG 1 .890 .902 rg/' \
> -e's/\.627 \.322 \.176 RG \.627 \.322 \.176 rg/.184 .31 .31 RG .184 .31 .31 rg/'
\end{verbatim}
\endgroup
\par\vspace{1.5pt}
\centering
\includegraphics[trim={0 7mm 0 75mm}, clip, width=\linewidth, page=406]{images/screen2-texts-replaced-color-compressed}
\par\vspace{-1.5pt}
\caption{\acro{PDF} dokument s řeckým a anglickým textem jsem dekomprimoval nástrojem \acro{PDF}tk a následně jsem ho upravil \acro{POSIX}ovým editorem \texttt{sed} tak, aby byl anglický text byl na červeném poli a spolu s řeckým textem tak tvořil horní dvě barvy české trikolory.}
\label{fig:screen2-texts-greek-and-english-updated}
\end{subfigure}
\par\vspace{10pt}
\begin{subfigure}{\linewidth}
\begingroup
\footnotesize
\begin{verbatim}
$ parallel -- convert -density 600 {} -resize '1920x1200!' sRGB {.}.png ::: *.pdf
$ ROZMAZEJ='-resize 50% -filter Gaussian -define filter:sigma=12 -resize 200%'
$ convert pozadí.png \( \( zářící-text.png $ROZMAZEJ \) základní-text.png \
>                       -composite \) -composite výstup.png
\end{verbatim}
\endgroup
\par\vspace{1.5pt}
\centering
\begin{tikzpicture}[every node/.style={inner sep=0pt, outer sep=0pt}]
\node at (0, 0) {\includegraphics[trim={0 85mm 0 6.8mm}, clip, width=\linewidth, page=406]{images/screen2-czech/fantasia-background}};
\node at (0, 0) {\includegraphics[trim={0 286.9mm 0 22.9mm}, clip, width=\linewidth]{images/screen2-czech/fantasia-effects-406}};
\node at (0, 0) {\includegraphics[trim={0 85mm 0 6.8mm}, clip, width=\linewidth, page=406]{images/screen2-czech/fantasia-text}};
\end{tikzpicture}
\par\vspace{-1.5pt}
\caption{Pro spodní část panelu jsem pomocí \TeX u vytvořil samostatný dokument s českým ekumenickým překladem Janovy Apokalypsy. Pro efekt záře jsem vytvořil tři verze dokumentu s pozadím, zářícím textem a základním textem. Dokumenty jsem rasterizoval na bitmapy, zářící text jsem rozmazal a bitmapy jsem překryl, vše nástrojem ImageMagick.}
\label{fig:screen2-texts-czech}
\end{subfigure}
\caption{Příprava panelu s trojjazyčným textem Zjevení Janova}
\label{fig:screen2-texts}
\end{figure}